% !PRINTED 452
\weberpage{452}
%<<CONT
adjungirt, den K\"orper der complexen Zahlen $x+yi$, mit rationalen $x, y$, den wir $\mathfrak{J}$ nennen wollen. Durch Adjunction einer Variablen $u$ zum K\"orper $\mathfrak{R}$ erh\"alt man den K\"orper, der aus allen ganzen und gebrochenen rationalen Functionen von $u$ mit rationalen Co\"efficienten besteht, durch Adjunction von $u$ und $i$ zu $\mathfrak{R}$ den K\"orper der rationalen Functionen von $u$, deren Co\"efficienten Zahlen in $\mathfrak{J}$ sind u.\,s.\,f.

Es sei schliesslich noch bemerkt, dass ein \gesperrt{K\"orper} von Gr\"ossen vielfach auch, nach \gesperrt{Kronecker}'s Vorgang, ein \gesperrt{Rationalit\"atsbereich} genannt wird. Es ist bisweilen n\"utzlich, den Ausdruck Rationalit\"atsbereich f\"ur einen K\"orper dann zu brauchen, wenn damit ausgedr\"uckt werden soll, dass bei einer bestimmten Aufgabe die Gr\"ossen dieses K\"orpers als bekannt oder rational betrachtet werden sollen.

\begin{center}
\S.\,141.

\gesperrt{Functionen in einem K\"orper.}
\end{center}

Wir haben oft Veranlassung, ganze rationale Functionen einer Ver\"anderlichen $x$
\weq{1}{f(x) = a_0 x^m + a_1 x^{m-1} + \cdots + a_{m-1} x + a_m}
zu betrachten, deren Co\"efficienten einem bestimmten K\"orper $\Omega$ angeh\"oren, und die wir in $\Omega$ \gesperrt{enthaltene Functionen}, oder kurz \gesperrt{Functionen in} $\Omega$ nennen wollen. Sie sind nicht zu verwechseln mit Functionen, die etwa dem K\"orper $\Omega$ \gesperrt{angeh\"oren}, wenn $\Omega$ ein Functionenk\"orper ist.

Im dritten Abschnitt haben wir gesehen, dass, wenn die Co\"efficienten $a_0, a_1 \ldots a_m$ bestimmte Zahlenwerthe haben, $m$ Wurzeln $\alpha_1, \alpha_2 \ldots \alpha_m$ von $f(x)$ existiren, die gleichfalls bestimmte Zahlenwerthe haben, unter denen aber mehrere einander gleiche vorkommen, wenn die Discriminante $\Delta(a)$ der Function (1) verschwindet. Sind die Co\"efficienten $a$ von $f(x)$ unabh\"angige Variable, so giebt es f\"ur jedes bestimmte Werthsystem der $a$ ein bestimmtes Werthsystem der $\alpha$ und die $\alpha$ \"andern sich stetig mit den $a$ (\S.\,40). So lange wir also die Ver\"anderung der $a$ auf ein hinl\"anglich enges Gebiet beschr\"anken, in dem solche Werthe nicht vorkommen, f\"ur welche die Discriminante $\Delta(a)$ verschwindet, so sind die $m$ Gr\"ossen $\alpha$ als Functionen der $a$ vollst\"andig bestimmt und von einander unterschieden, wenn der Werth einer jeden von ihnen f\"ur irgend ein bestimmtes Werthsystem der $a$ gegeben ist.
%CONT>>
